\documentclass[a4paper,fontset=none]{ctexart}

\usepackage{anyfontsize}
\usepackage{amsmath,amssymb,mathrsfs,wasysym}
\allowdisplaybreaks
\usepackage[T1]{fontenc}
\usepackage{textcomp}
\usepackage{amsthm}
\usepackage[libertine,vvarbb]{newtxmath}
\usepackage[scr=rsfso,frak=euler,bb=ams]{mathalfa}
\usepackage{bm}
\usepackage{indentfirst}
\setlength{\parindent}{0em}
\usepackage{parskip}
\usepackage{geometry}
\geometry{top=3cm,bottom=3cm,left=2cm,right=2cm}
\usepackage{fontspec}
\setmainfont{Linux Libertine O}
\setsansfont{Linux Biolinum O}
\setmonofont{JetBrains Mono}
\setCJKmainfont{SimSun}[BoldFont=Noto Sans CJK SC Medium]

\begin{document}

    \title{\textbf{天体物理导论作业}}
    \author{1900011604\ \ 张植竣\ \ 北京大学}
    \date{2021年3月8日}

    \maketitle

    \section*{第一章}

    \begin{itemize}
        \item[1.]
        三角视差法：测量视差角$\pi$和基线长度$l$，则距离$D$满足：
        \[
        \sin\pi = \frac{l}{D} \approx \pi
        \]
        则计算出：
        \[
        D = \frac{l}{\sin\pi} = \frac{l}{\pi}
        \]

        HR图法：选定一些主序星，通过测量它们的颜色（光谱）确定温度$T$，再利用已有主序星在HR图上的分布确定光度$L$。再测量视亮度$F$，计算出距离：
        \[
        D = \sqrt{\frac{L}{4\pi F}}
        \]

        造父变星法：某些恒星在引力平衡态附近出现周期性的振荡，并且其周期$P$与光度$L$动力学相关，称为周光关系，造父变星是这类恒星的典型。测量其光变周期$P$可以计算出光度$L$，再测量视亮度$F$，用$L=4\pi D^2 F$计算距离$D$。

        超新星法：认为某类超新星（例如Ia型超新星）爆发时的光度是常值$L=L_0$，因为它们有共同的物理过程，再测量视亮度$F$，用$L=4\pi D^2 F$计算距离$D$。

        谱线宽度法（星系）：对漩涡星系使用Tully-Fisher关系，测量谱线宽度$\sigma$，则光度$L \propto \sigma^4$；或对椭圆星系使用Faber-Jackson关系，测量发光物质的速度弥散（谱线宽度）$\sigma$，则光度$L \propto \sigma^4$。再测量视亮度$F$，用$L=4\pi D^2 F$计算距离$D$。

        Hubble关系法：红移$z$较小（此时Hubble常数$H$随时间的变化可以忽略）时，利用Hubble关系$v=Hr$，通过红移$z$测量远处天体的退行速度$v=cz$，就可以计算出到观测者的距离：
        \[
        r = \frac{v}{H} = \frac{cz}{H}
        \]

        \item[2.]
        由于各辐射单元几乎同时光变，且几乎是各向同性地发光，该天体的尺度$L$即为最早观测到光变的位置至最晚观测到光变的位置的距离，故$L \approx \tau c$。

    \end{itemize}

    \section*{第二章}

    \begin{itemize}
        \item[1.]
        由题意：
        \[
        F = \frac{L}{4\pi d^2} = \frac{S\sigma T^4}{4\pi d^2} = 1.7\times10^{-13}\,\mathrm{erg\cdot cm^{-2}\cdot s^{-1}}
        \]
        则黑体辐射区的面积：
        \[
        S = \frac{4\pi d^2 F}{\sigma T^4} = 3.20\times10^{7}\,\mathrm{cm^2}
        \]

        \item[2.(1)]
        最低辐射为回旋辐射，其基频为：
        \[
        \omega_0 = \frac{eB}{mc} = 1.759\times10^3\,\mathrm{rad/s}
        \]
        即
        \[
        \nu_0 = \frac{\omega_0}{2\pi} = 2.80\times10^2\,\mathrm{Hz}
        \]

        \item[2.(2)]
        由于Lorentz因子$\gamma = 100 \gg 1$，电子与磁场作用以同步辐射为主，辐射谱最强的频率为：
        \[
        \nu_\mathrm{m}^s \approx 0.45\gamma^3\nu_0 = 1.26\times10^6\,\mathrm{Hz}
        \]

        \item[2.(3)]
        取可见光中心波长$\lambda = 555\,\mathrm{nm}$，则有：
        \[
        0.45\gamma^3\nu_0 = \frac{c}{\lambda}
        \]
        则Lorentz因子为：
        \[
        \gamma = \frac{c\mu_0}{0.45\lambda} = 6.952\times10^{5}
        \]

        \item[4.]
        逆Compton散射的功率为：
        \[
        \frac{\mathrm{d}E}{\mathrm{d}t}
        = mc^2\frac{\mathrm{d}\gamma}{\mathrm{d}t}
        = \frac{32\pi}{9}r_e^2 c\rho\gamma^2
        \]
        整理得：
        \[
        \frac{9mc}{32\pi r_e^2\rho\gamma^2}\mathrm{d}\gamma = \mathrm{d}t
        \]
        两边积分得到：
        \[
        \int_{\gamma_0-\Delta\gamma}^{\gamma_0} \frac{9mc}{32\pi r_e^2\rho\gamma^2}\mathrm{d}\gamma = \int_0^\tau \mathrm{d}t
        \]
        即
        \[
        \frac{9mc}{32\pi r_e^2\rho(\gamma_0 - \Delta\gamma)}-\frac{9mc}{32\pi r_e^2\rho\gamma_0} = \tau
        \]
        代入$\rho = \dfrac{4\sigma}{c}T^4$得：
        \[
        \tau = \frac{9mc^2}{128\pi r_e^2 \sigma T^4(\gamma_0 - \Delta\gamma)}-\frac{9mc^2}{128\pi r_e^2 \sigma T^4\gamma_0}
        \]
        考虑Lorentz因子的变化$\Delta\gamma = \dfrac{\gamma_0}{2} = 50$，则$\tau = 5.024\times10^{17}\,\mathrm{s}$。

        此时电子的速度仍然十分接近于光速，行走的距离为：
        \[
        d \approx c\tau = 1.506\times10^{28}\,\mathrm{cm} = 4.881\times10^{9}\,\mathrm{pc}
        \]

        \item[4.]
        两个电子碰撞时，因为两个粒子的质量和电荷量都相等，在质心系中：
        \[
        \bm{v}_1 + \bm{v_2} = 0,\qquad \bm{x}_1 + \bm{x_2} = 0
        \]
        则系统的电偶极矩为零：
        \[
        \bm{p} = e\bm{x}_1 + e\bm{x}_2 = 0
        \]
        故不发生电偶极辐射。
    \end{itemize}

    \section*{第三章}

    \begin{itemize}

        \item[2.]
        设光球层的粒子数密度为$n$，由题意得：
        \[
        \frac{B_1^2}{8\pi} + nkT_1 = \frac{B_2^2}{8\pi} + nkT_2
        \]
        即：
        \[
        n = \frac{B_2^2 - B_1^2}{8\pi k(T_1 - T_2)}
        \]
        其中$B_1 = 50\,\mathrm{G}$，$B_2 = 1000\,\mathrm{G}$，$T_1 - T_2 = 2000\,\mathrm{K}$，计算得$n = 1.437\times10^{17}\,\mathrm{cm^{-3}}$，则物质密度为：
        \[
        \rho = m_{\mathrm{H}}n = 2.406\times10^{-7}\,\mathrm{g/cm^3}
        \]

        \item[3.]
        由于磁冻结现象，恒星塌缩过程中磁通量$\varPhi$守恒，而$\varPhi = \bm{B}\cdot\bm{S} \propto BR^2$，故磁场与半径的关系为$B \propto R^{-2}$。

    \end{itemize}
\end{document}
